\subsection{正合列}

回顾下面的定义（《代数学》第二章 §1 第 4 目）：

定义 1. 设 A 是环，F, G, H 是三个右（相应地左）A-模，f 是 F 到 G 的同态，g 是 G 到 H 的同态。称有序偶 \((f, g)\) 为一个正合列，如果 \(\bar{g}^{-1}(0) = f(F)\) ，即如果 g 的核等于 f 的像。

图

\[
\mathrm{F} \xrightarrow {f} \mathrm{G} \xrightarrow {g} \mathrm{H}\tag{4}
\]

也称为一个正合列。

类似地，考虑一个由四个模和三个同态组成的图：

\[
\mathrm{E} \xrightarrow {f} \mathrm{F} \xrightarrow {g} \mathrm{G} \xrightarrow {h} \mathrm{H}.\tag{5}
\]

若图 \(E \xrightarrow{f} F \xrightarrow{g} G\) 是一个正合列，则称这个图在 F 处正合；若 \(F \xrightarrow{g} G \xrightarrow{h} H\) 是一个正合列，则称它在 G 处正合。若 (5) 在 F 和 G 处都正合，则称它是正合的，或也称它为一个正合列。任意多项的正合列可类似地定义。

回顾下面的结果（出处同上），其中 E, F, G 表示右（相应地左）

A-模，箭头表示同态，而 0 表示仅由其单位元组成的模：

\begin{enumerate}
\def\labelenumi{(\alph{enumi})}
\item
  说 \(0 \to \mathbf{E} \xrightarrow{f} \mathbf{F}\) 是一个正合列，等价于说 \(f\) 是单射。
\item
  说 \(\mathbf{E} \xrightarrow{f} \mathbf{F} \to 0\) 是一个正合列，等价于说 \(f\) 是满射。
\item
  说 \(0 \to \mathbf{E} \xrightarrow{f} \mathbf{F} \to 0\) 是一个正合列，等价于说 \(f\) 是双射，即 \(f\) 是 \(\mathbf{E}\) 到 \(\mathbf{F}\) 上的同构。
\item
  若 F 是 E 的子模，i 表示 F 到 E 中的典范单射，p 表示 E 到 E/F 上的典范满射，则图
\end{enumerate}

\[
0 \longrightarrow \mathrm{F} \stackrel {i} {\longrightarrow} \mathrm{E} \stackrel {p} {\longrightarrow} \mathrm{E} / \mathrm{F} \longrightarrow 0\tag{6}
\]

是一个正合列。

\begin{enumerate}
\def\labelenumi{(\alph{enumi})}
\setcounter{enumi}{4}
\tightlist
\item
  若 \(f: \mathbf{E} \to \mathbf{F}\) 是同态，则图
\end{enumerate}

\[
0 \longrightarrow \overline {{f}} ^ {- 1} (0) \stackrel {i} {\longrightarrow} \mathrm{E} \stackrel {f} {\longrightarrow} \mathrm{F} \stackrel {p} {\longrightarrow} \mathrm{F} / f (\mathrm{E}) \longrightarrow 0\tag{7}
\]

（其中 \(i\) 是 \(\overline{f}(0)\) 到 E 中的典范单射， \(p\) 是 \(\mathbf{F} / f(\mathbf{E})\) 的典范投影）是一个正合列。

\begin{enumerate}
\def\labelenumi{(\alph{enumi})}
\setcounter{enumi}{5}
\tightlist
\item
  为使图
\end{enumerate}

\[
\mathrm{E} \xrightarrow {f} \mathrm{F} \xrightarrow {g} \mathrm{G}\tag{8}
\]

成为一个正合列，必要且充分的是：存在模 S, T 以及同态 \(a: \mathrm{E} \to \mathrm{S}\) 、 \(b: \mathrm{S} \to \mathrm{F}\) 、 \(c: \mathrm{F} \to \mathrm{T}\) 和 \(d: \mathrm{T} \to \mathrm{G}\) ，使得 \(f = b \circ a\) 、 \(g = d \circ c\) ，且三个序列

\[
\begin{array}{c} \mathrm{E} \xrightarrow {a} \mathrm{S} \longrightarrow 0 \\ 0 \longrightarrow \mathrm{S} \xrightarrow {b} \mathrm{F} \xrightarrow {c} \mathrm{T} \longrightarrow 0 \\ 0 \longrightarrow \mathrm{T} \xrightarrow {d} \mathrm{G} \end{array}\tag{9}
\]

都是正合的。

最后回顾，若 \(f: \mathbf{E} \to \mathbf{F}\) 是 A-模同态，我们令 \(\operatorname{Ker}(f) = \bar{f}^{-1}(0)\) ， \(\operatorname{Im}(f) = f(\mathbf{E})\) ， \(\operatorname{Coim}(f) = \mathbf{E} / \bar{f}^{-1}(0)\) 和 \(\operatorname{Coker}(f) = \mathbf{F} / f(\mathbf{E})\) 。用这一记号，可以在 (9) 中取 \(\mathbf{S} = \operatorname{Im}(f) = \operatorname{Ker}(g)\) 和 \(\mathbf{T} = \operatorname{Im}(g)\) （典范同构于 \(\operatorname{Coker}(f)\) ）。

